\chapter{Il livello applicativo: architetture e servizi}

Cominciamo la discesa lungo la pila partendo, come vuole l'approccio top-down, dallo strato più
alto: quello che l'utente vede.

\begin{center}
\begin{tikzpicture}[node distance=0pt]
  \node[lay app, minimum width=3.4cm, line width=1.5pt, draw=netred] (a) {\bfseries Applicazione};
  \node[lay tra, minimum width=3.4cm, below=of a] (t) {Trasporto};
  \node[lay net, minimum width=3.4cm, below=of t] (n) {Rete};
  \node[lay link, minimum width=3.4cm, below=of n] (l) {Collegamento};
  \node[lay phy, minimum width=3.4cm, below=of l] (p) {Fisico};
  \draw[-{Stealth}, thick, netred] ([xshift=0.3cm]a.east) -- ++(1,0) node[right, lbl, text=netred] {siamo qui};
\end{tikzpicture}
\end{center}

Abbiamo l'intera pila; ora apriamo lo strato in cima.
\begin{itemize}
    \item Un'applicazione \textbf{non sposta dati}: li consegna allo strato sotto e li ritrova all'altro capo.
    \item Il trasporto le offre (se lo chiede) un \textbf{flusso di dati affidabile}: niente perso, niente fuori ordine.
    \item Due applicazioni su due host si comportano \textbf{come se parlassero direttamente}, purché seguano lo
          stesso protocollo.
\end{itemize}

\info{Che cosa resta da fare al programmatore?}{
    Tutto ciò che sta \textbf{sopra} quel flusso: \textbf{quali messaggi} esistono, \textbf{che cosa significano}
    e \textbf{chi parla per primo}. È esattamente ciò che definisce un protocollo applicativo.
}

% ══════════════════════════════════════════════════════════════
\section{Che cos'è un'applicazione di rete}
% ══════════════════════════════════════════════════════════════

\dfn{Applicazione di rete}{
    Un'\textbf{applicazione di rete} è composta da \textbf{più programmi} che girano su
    \textbf{end system diversi} (host) e che \textbf{comunicano tra loro attraverso la rete}.
}

\begin{esempio}{Un'applicazione Web}
Un'applicazione Web è fatta di due programmi distinti:
\begin{itemize}
    \item il \textbf{browser}, in esecuzione sull'host dell'utente (desktop, laptop, tablet, smartphone);
    \item il \textbf{server Web}, in esecuzione sull'host che ospita il sito.
\end{itemize}
Possono essere scritti in linguaggi diversi e girare su sistemi operativi diversi: l'unica cosa che
devono condividere è il \textbf{protocollo} (HTTP).
\end{esempio}

Qualche applicazione di rete tipica: social network, Web, messaggistica, posta elettronica, giochi
multiutente, streaming video (YouTube, Netflix), condivisione di file P2P, voice over IP,
videoconferenza, ricerca su Internet, login remoto.

\warning{Il Web è tecnicamente un'applicazione}{
    Nel linguaggio comune ``Internet'' e ``Web'' sono usati come sinonimi, ma il Web è
    \textbf{una delle applicazioni} che girano sopra Internet, non è Internet.
}

\subsection{Creare un'applicazione di rete}

\begin{figure}[H]
\centering
\begin{minipage}{0.4\linewidth}\centering
  \includegraphics[width=0.9\linewidth]{cap03/app_stack_endsystem.png}
\end{minipage}\hfill
\begin{minipage}{0.56\linewidth}
Creare un'applicazione di rete significa scrivere programmi che:
\begin{itemize}
  \item girano su \textbf{end system diversi}, magari scritti in linguaggi diversi o per sistemi
        operativi diversi;
  \item comunicano attraverso la rete (per esempio tramite \textbf{socket}) usando uno
        specifico \textbf{protocollo}.
\end{itemize}
Notate nella figura che la pila completa (applicazione, trasporto, rete, collegamento, fisico) c'è
\textbf{solo sugli host} ai bordi della rete.
\end{minipage}
\caption{Le applicazioni girano solo sugli end system.}
\end{figure}

\info{Non si scrive software per il nucleo della rete}{
    Non c'è bisogno di scrivere software per i dispositivi del \emph{network core} (router, switch):
    \begin{itemize}
      \item i dispositivi di rete \textbf{non eseguono applicazioni utente};
      \item esistono librerie specifiche (i \textbf{socket}) che implementano le funzionalità di rete.
    \end{itemize}
    È una conseguenza diretta del modello a strati: l'applicazione vede soltanto il servizio che le offre
    il trasporto, e ignora come sia realizzato. Router e switch non eseguono applicazioni utente, e non vi
    lascerebbero installarne una. Questa restrizione \textbf{non è un limite}: è proprio il motivo per cui
    nuove applicazioni nascono e si diffondono così in fretta, senza toccare l'infrastruttura.
}

% ══════════════════════════════════════════════════════════════
\section{Architetture delle applicazioni di rete}
% ══════════════════════════════════════════════════════════════

\dfn{Architettura applicativa}{
    L'\textbf{architettura} di un'applicazione di rete è progettata dallo \textbf{sviluppatore}
    dell'applicazione e determina come l'applicazione è strutturata sui vari end system.
    Le due architetture fondamentali sono \textbf{client-server} e \textbf{peer-to-peer}.
}

\warning{Architettura applicativa e architettura di rete}{
    Non confondete l'\emph{architettura dell'applicazione} (decisa dallo sviluppatore: client-server o P2P)
    con l'\emph{architettura della rete} (la pila a cinque strati di Internet, che è fissata e offre
    all'applicazione un insieme di servizi).
}

\begin{figure}[H]
\centering
\begin{tikzpicture}[every node/.style={font=\footnotesize\sffamily}]
  \begin{scope}
    \node[font=\small\bfseries\sffamily] at (0,2.5) {Client-server};
    \node (srv) at (0,0) {\icon[0.8cm]{server}};
    \node[below=-1pt of srv, text=netred] {server sempre attivo};
    \foreach \a/\ic in {90/laptop, 30/pc, -30/smartphone, -90/laptop, 210/pc, 150/smartphone} {
      \node (c\a) at (\a:2.1) {\icon[0.65cm]{\ic}};
      \draw[msg, netblue, {Stealth}-{Stealth}] (c\a) -- (srv);
    }
    \node[lbl] at (0,-3.1) {i client parlano solo con il server};
  \end{scope}
  \begin{scope}[xshift=8cm]
    \node[font=\small\bfseries\sffamily] at (0,2.5) {Peer-to-peer};
    \foreach \a/\ic [count=\i] in {90/laptop, 30/pc, -30/laptop, -90/pc, 210/laptop, 150/pc} {
      \node (p\i) at (\a:1.9) {\icon[0.65cm]{\ic}};
    }
    \foreach \x/\y in {1/2, 1/4, 2/3, 2/5, 3/6, 4/5, 5/6, 6/1, 3/4} \draw[msg, netgreen, {Stealth}-{Stealth}] (p\x) -- (p\y);
    \node[lbl] at (0,-3.1) {i peer comunicano direttamente tra loro};
  \end{scope}
\end{tikzpicture}
\caption{Le due architetture applicative fondamentali.}
\end{figure}

\subsection{Architettura client-server}

\dfn{Architettura client-server}{
    I nodi sono \textbf{eterogenei}: esiste un host \textbf{sempre attivo} (il \emph{server}) che fornisce
    servizi in risposta alle richieste di molti altri host (i \emph{client}).
}

\begin{itemize}
    \item I client \textbf{non comunicano direttamente} tra loro: due browser non si scambiano mai nulla.
    \item Il server ha un \textbf{indirizzo fisso e ben noto} (un nome o un indirizzo IP): un client può
          sempre contattarlo inviandogli un pacchetto di richiesta.
    \item \textbf{Il client parla per primo}: il server aspetta di essere contattato, e risponde.
    \item Spesso servono \textbf{molti server} per reggere tutte le richieste; il client di solito non se
          ne accorge e li percepisce come un server unico.
\end{itemize}

I server multipli possono essere organizzati in tre modi:

\begin{itemize}
    \item \textbf{raggruppati in data center}, con un numero enorme di server in un unico luogo: vanno
          alimentati, raffreddati, mantenuti e ben collegati;
    \item \textbf{sparsi come server distribuiti} in tutto il mondo: vanno interconnessi e coordinati;
    \item organizzati in \textbf{data center distribuiti}, combinando le due soluzioni: è ciò che fanno davvero
          i grandi fornitori.
\end{itemize}

E il fornitore continua a pagare, mese dopo mese, l'interconnessione e la banda che esce da quegli edifici.

\warning{La conseguenza del client-server}{
    Tutto ciò che l'applicazione fa passa per il server, che è quindi insieme il \textbf{limite alla
    scalabilità} (quanti client si possono servire) e il \textbf{punto di guasto unico} (\emph{single point of
    failure}): se il server cade, il servizio non esiste più per nessuno.
}

\begin{esempio}{Il Web come applicazione client-server}
C'è un server Web sempre attivo che riceve richieste dai browser degli utenti. Quando riceve la
richiesta di un oggetto (una pagina, un'immagine), risponde inviando l'oggetto richiesto. Il server
è raggiungibile da tutti gli host, perché il suo indirizzo è noto, ed è acceso quando lo si interroga.
Richiesta, risposta, e un server che c'era già prima che qualcuno chiedesse: è tutta qui la forma del
client-server.
\end{esempio}

\subsection{Architettura peer-to-peer (P2P)}

\dfn{Architettura peer-to-peer}{
    I nodi sono (idealmente) \textbf{omogenei}: l'applicazione sfrutta la \textbf{comunicazione diretta}
    tra coppie di host collegati in modo \textbf{intermittente}, detti \emph{peer}. I peer non sono di
    proprietà del fornitore del servizio, ma sono computer controllati dagli utenti.
}

\begin{itemize}
    \item È usata spesso per applicazioni ad \textbf{alta intensità di traffico} (condivisione di file).
    \item Non c'è un server unico con indirizzo fisso: ogni peer ha il proprio indirizzo, che può anche cambiare.
    \item È \textbf{auto-scalabile} (\emph{self-scalability}): ogni peer genera carico chiedendo un file, ma
          aggiunge anche \textbf{capacità} servendolo ad altri. La capacità cresce insieme alla domanda: nessun
          progetto client-server ha questa proprietà.
    \item È \textbf{economica}: non servono infrastrutture né grande banda da parte del fornitore.
    \item Ha però problemi di \textbf{sicurezza}, \textbf{prestazioni} e \textbf{affidabilità}: nessuno possiede
          le macchine coinvolte.
\end{itemize}

\info{Il P2P puro è raro: le architetture ibride}{
    Le applicazioni P2P pure sono rare: la maggior parte ha un'architettura \textbf{ibrida}. In molte
    applicazioni di messaggistica e di telefonia su Internet un \textbf{server} tiene traccia degli
    indirizzi IP degli utenti, ma i messaggi da utente a utente viaggiano \textbf{direttamente} tra gli
    host, senza passare per server intermedi. Nella condivisione di file il server può anche tenere
    l'elenco dei file disponibili, per velocizzare la ricerca (così facevano Napster e, in parte,
    BitTorrent con i tracker).
}

\begin{figure}[H]
\centering
\begin{tikzpicture}[every node/.style={font=\footnotesize\sffamily}]
  \node (S) at (0,2.2) {\icon[0.7cm]{server}};
  \node[right=2pt of S, lbl, text width=3.2cm, align=left] {server di indicizzazione: ``Bob è all'indirizzo X''};
  \node (A) at (-3.5,0) {\icon[1cm]{laptop}}; \node[below=-2pt of A] {Alice};
  \node (B) at (3.5,0) {\icon[0.9cm]{pc}}; \node[below=-2pt of B] {Bob};
  \draw[msg, dashed, netblue] (A) -- node[above, sloped] {1. dov'è Bob?} (S);
  \draw[msg, dashed, netblue] (B) -- node[above, sloped] {0. registrazione} (S);
  \draw[msg, very thick, netgreen, {Stealth}-{Stealth}] (A) -- node[below=2pt] {2. messaggi diretti, peer-to-peer} (B);
\end{tikzpicture}
\caption{Un'architettura ibrida: il server serve solo a trovarsi, poi i peer parlano direttamente.}
\end{figure}

% ══════════════════════════════════════════════════════════════
\section{Comunicazione tra processi}
% ══════════════════════════════════════════════════════════════

Come abbiamo già detto, chi comunica \textbf{non è un programma e non è una macchina}, ma un
\textbf{processo}: un programma in esecuzione, su macchine diverse e potenzialmente con sistemi operativi diversi.

\begin{itemize}
    \item Nel Web il processo del browser scambia messaggi con il processo del server Web.
    \item Nella condivisione di file P2P un file passa dal processo di un peer al processo di un altro peer.
\end{itemize}

\dfn{Processo client e processo server}{
    In ogni coppia di processi che comunicano:
    \begin{itemize}
      \item il \textbf{processo client} è quello che \textbf{inizia} la comunicazione;
      \item il \textbf{processo server} è quello che \textbf{attende} di essere contattato.
    \end{itemize}
}

Nel Web il browser è il client e il server Web è il server. Nel P2P il peer che scarica fa da
client e quello che carica fa da server: in un'applicazione P2P i processi possono
\textbf{cambiare ruolo}, e lo stesso peer può essere client su una connessione e server su un'altra.

\subsection{I socket}

La maggior parte delle applicazioni è fatta di coppie di processi che si mandano messaggi
attraverso la rete sottostante. Il punto in cui un processo ``consegna'' i messaggi alla rete ha un
nome preciso.

\dfn{Socket}{
    Seguendo il modello a strati, un \textbf{socket} è l'\textbf{interfaccia software tra il livello
    applicativo e il livello di trasporto}: è la ``porta'' attraverso cui un processo invia e riceve
    messaggi sulla rete. È anche l'\textbf{API} (\emph{Application Programming Interface}) tra
    l'applicazione e la rete.
}

\begin{figure}[H]
    \centering
    \includegraphics[width=0.72\linewidth]{cap03/processi_socket.png}
    \caption{Il processo è controllato dallo sviluppatore; ciò che sta sotto il socket (TCP con i suoi buffer e le sue variabili) è controllato dal sistema operativo.}
\end{figure}

\begin{esempio}{I socket come cassette postali}
\begin{center}
\includegraphics[width=0.42\linewidth]{cap03/socket_cassetta.png}
\end{center}
\begin{itemize}
    \item Un processo che vuole mandare un messaggio a un processo su un altro host \textbf{lo mette nella
          cassetta postale} (il socket).
    \item Il processo mittente dà per scontato che, dall'altra parte della sua porta, ci sia
          un'\textbf{infrastruttura di trasporto} che porterà il messaggio fino alla cassetta del destinatario.
    \item Quando il messaggio arriva all'host di destinazione, passa attraverso la cassetta (il socket) del
          processo ricevente.
\end{itemize}
Chi imbuca una lettera non deve sapere come funziona il sistema postale.
\end{esempio}

\begin{figure}[H]
\centering
\begin{tikzpicture}[every node/.style={font=\footnotesize\sffamily}]
  \newcommand{\casa}[2]{%
    \begin{scope}[shift={(#1)}]
      \draw[thick, fill=lapp!40] (0,0) rectangle (1.8,1.3);
      \draw[thick, fill=netred!50] (-0.2,1.3) -- (0.9,2.1) -- (2.0,1.3) -- cycle;
      \draw[thick, fill=ltra] (1.2,0) rectangle (1.65,0.8);
      \fill (1.3,0.4) circle (1.2pt);
      \node[font=\scriptsize\bfseries\sffamily] at (0.55,0.75) {#2};
    \end{scope}}
  \casa{0,0}{processo}
  \casa{10,0}{processo}
  \node[text=netblue] at (1.45,-0.3) {socket = porta};
  \node[text=netblue] at (11.45,-0.3) {socket = porta};
  \node[netcloud, minimum width=4.5cm, minimum height=1.6cm] (net) at (6.2,0.55) {infrastruttura di trasporto};
  \draw[-{Stealth}, very thick, netred] (1.7,0.45) -- node[above, lbl] {messaggio} (net.west);
  \draw[-{Stealth}, very thick, netred] (net.east) -- (11.25,0.45);
  \node[lbl, text width=3.2cm] at (1,-1.1) {spinge il messaggio fuori dalla propria porta};
  \node[lbl, text width=3.2cm] at (11,-1.1) {lo riceve dalla propria porta, e solo allora agisce};
\end{tikzpicture}
\caption{Un processo è una casa e il suo socket è la porta: il mittente spinge fuori il messaggio e dà per scontato che qualcosa, fuori, lo porti a destinazione.}
\end{figure}

\warning{Cosa controlla lo sviluppatore}{
    Lo sviluppatore ha il \textbf{controllo completo} di tutto ciò che sta sul lato applicativo del
    socket, ma \textbf{pochissimo controllo} sul lato del trasporto. Può soltanto:
    \begin{enumerate}
      \item \textbf{scegliere il protocollo di trasporto} (TCP o UDP);
      \item eventualmente impostare qualche parametro del trasporto (dimensione massima dei buffer, dei
            segmenti, \dots).
    \end{enumerate}
    Scelto il protocollo, l'applicazione viene costruita dando per scontati i servizi che quel protocollo offre.
    La scelta si fa \textbf{una volta, all'inizio}, e tutto il resto ne discende: è una \textbf{decisione di
    progetto}, non un dettaglio di configurazione.
}

\subsection{Indirizzare i processi: indirizzo e identificatore}

Riprendendo l'analogia postale, un processo ha bisogno di un \textbf{indirizzo} per ricevere
messaggi. Ma su un host girano contemporaneamente molte applicazioni di rete: per individuare il
processo destinatario servono \textbf{due} informazioni.

\begin{enumerate}
    \item Un \textbf{indirizzo dell'host}, per trovare la macchina giusta sulla rete.
    \item Un \textbf{identificatore del processo}, per trovare il processo giusto dentro quella macchina.
\end{enumerate}

\begin{figure}[H]
\centering
\begin{tikzpicture}[every node/.style={font=\footnotesize\sffamily}]
  \node[draw=netgreen!70!black, thick, rounded corners, minimum width=3.6cm, minimum height=3cm] (h1) at (0,0) {};
  \node[above=1pt of h1, text=netgreen!60!black] {telefono \texttt{93.45.7.12}};
  \node[ellipse, draw, fill=lapp!50, minimum width=1.9cm] (tg) at (-0.4,0.7) {Telegram};
  \node[ellipse, draw, fill=lapp!50, minimum width=1.9cm] (br) at (-0.4,-0.7) {Opera};
  \node[field, fill=ltra, minimum width=0.4cm] (p1) at (1.75,0.7) {};
  \node[field, fill=ltra, minimum width=0.4cm] (p2) at (1.75,-0.7) {};
  \draw (tg) -- (p1); \draw (br) -- (p2);
  \node[above=-1pt of p1, text=netblue, font=\scriptsize\ttfamily] {51001}; \node[below=-1pt of p2, text=netblue, font=\scriptsize\ttfamily] {51000};
  \node[netcloud, minimum width=3.4cm, minimum height=2cm] (net) at (5.6,0) {rete};
  \node[draw=netgreen!70!black, thick, rounded corners, minimum width=2.8cm, minimum height=1.2cm] (h2) at (10.6,1.4) {};
  \node[above=1pt of h2, text=netgreen!60!black] {un altro telefono};
  \node[ellipse, draw, fill=lapp!50] (tg2) at (10.8,1.4) {Telegram};
  \node[field, fill=ltra, minimum width=0.4cm] (q1) at (9.15,1.4) {};
  \node[above=-1pt of q1, text=netblue, font=\scriptsize\ttfamily] {51002};
  \draw (tg2) -- (q1);
  \node[draw=netgreen!70!black, thick, rounded corners, minimum width=2.8cm, minimum height=1.2cm] (ws) at (10.6,-1.4) {};
  \node[below=1pt of ws, text=netgreen!60!black] {server Web UniNA};
  \node[ellipse, draw, fill=lapp!50] (ws2) at (10.8,-1.4) {server Web};
  \node[field, fill=ltra, minimum width=0.4cm] (q2) at (9.15,-1.4) {};
  \node[below=-1pt of q2, text=netblue, font=\scriptsize\ttfamily] {80};
  \draw (ws2) -- (q2);
  \draw[netred, thick, dashed, -{Stealth}] (p1) .. controls (4,1.7) and (7.5,1.9) .. (q1);
  \draw[netblue, thick, dashed, -{Stealth}] (p2) .. controls (4,-1.7) and (7.5,-1.9) .. (q2);
  \node[draw=primary, fill=primary!10, rounded corners, text width=7cm, align=center] at (5.6,-3.0) {\textbf{indirizzo}: livello di \textbf{rete}, trova l'host\\\textbf{porta}: livello di \textbf{trasporto}, trova il processo};
\end{tikzpicture}
\caption{Un telefono con la messaggistica e il browser aperti, che parlano con due destinazioni diverse nello stesso momento. Entrambi i flussi partono dallo stesso indirizzo: a tenerli separati, a entrambi i capi, è la porta.}
\end{figure}

\dfn{Indirizzo IP e numero di porta}{
    In Internet:
    \begin{itemize}
      \item un \textbf{host} è identificato da un \textbf{indirizzo IP} (in IPv4, una quantità a 32 bit,
            scritta come quattro numeri tra 0 e 255, per esempio \texttt{143.225.161.30});
      \item un \textbf{processo} è identificato da un \textbf{numero di porta} (16 bit).
    \end{itemize}
    L'indirizzo lavora a livello di \textbf{rete} (instradamento), l'identificatore a livello di \textbf{trasporto}.
}

Le applicazioni più diffuse hanno \textbf{numeri di porta assegnati per convenzione}: un server Web
ascolta sulla porta \textbf{80} (443 per HTTPS), un server di posta SMTP sulla porta \textbf{25}. L'elenco
ufficiale delle porte ben note è mantenuto dalla \textbf{IANA} (\emph{Internet Assigned Numbers Authority}).
È questa convenzione che permette a un client di contattare un servizio mai visto prima: senza, ogni client
dovrebbe sapere in anticipo ``a quale porta bussare''.

\subsection{Un indirizzo appartiene a un'interfaccia}

Un portatile è collegato contemporaneamente via cavo Ethernet e via Wi-Fi: quanti indirizzi ha? \textbf{Due},
uno per interfaccia. Un indirizzo IP non dà il nome a una macchina: dà il nome a \textbf{un punto di
collegamento a una rete}.

\begin{figure}[H]
\centering
\begin{tikzpicture}[every node/.style={font=\footnotesize\sffamily}]
  \node (lap) at (0,0) {\icon[1.3cm]{laptop}};
  \node[left=2pt of lap] {un portatile};
  \node[draw=netblue, fill=netblue!8, rounded corners, font=\scriptsize\ttfamily] (e) at (2.4,0.9) {eth0: 192.168.1.20};
  \node[draw=netgreen!70!black, fill=netgreen!8, rounded corners, font=\scriptsize\ttfamily] (w) at (2.4,-0.9) {wlan0: 192.168.1.35};
  \draw[wired] (lap) -- (e.west);
  \draw[wireless] (lap) -- (w.west);
  \node (sw) at (6,0.9) {\icon[0.9cm]{switch}};
  \node (ap) at (6,-0.9) {\icon[0.8cm]{accesspoint}};
  \draw[wired] (e) -- (sw); \draw[wireless] (w) -- (ap);
  \node (r) at (8.5,0) {\icon[1cm]{router}};
  \draw[wired] (sw) -- (r); \draw[wired] (ap) -- (r);
  \draw[netred, very thick] (4.1,1.15) -- (4.5,0.65); \draw[netred, very thick] (4.1,0.65) -- (4.5,1.15);
  \node[text=netred, lbl, anchor=west] at (4.6,1.45) {cavo staccato};
  \node[lbl, text width=4.2cm, anchor=west] at (9.3,0) {per il livello di rete sono \textbf{due host distinti}, raggiunti per due cammini diversi};
\end{tikzpicture}
\caption{Staccando il cavo, l'indirizzo di \texttt{eth0} smette di rispondere, mentre quello di \texttt{wlan0} continua a funzionare.}
\end{figure}

Ci torneremo quando gli indirizzi smetteranno di appartenere alle macchine e cominceranno ad appartenere
alle \textbf{reti} (capitolo sull'indirizzamento IP).

\subsection{IPv4 e l'indirizzo che vuol dire “qui”}

\dfn{Indirizzo IPv4 e localhost}{
    Un indirizzo \textbf{IPv4} è una quantità a \textbf{32 bit}, scritta come quattro numeri decimali tra 0 e 255
    separati da punti, uno per byte. L'indirizzo \texttt{127.0.0.1} (\texttt{localhost}) è speciale: non indica
    una macchina remota ma \textbf{questa} macchina. Due processi locali comunicano attraverso di esso
    esattamente come farebbero attraverso la rete, ma \textbf{nulla lascia l'host}.
}

\begin{figure}[H]
\centering
\begin{tikzpicture}[every node/.style={font=\footnotesize\sffamily}]
  \node[draw=netgreen!70!black, thick, rounded corners, minimum width=7cm, minimum height=2.4cm] (h) at (0,0) {};
  \node[anchor=north east, text=netgreen!60!black] at (h.north east) {host};
  \node[ellipse, draw, fill=lapp!50, minimum width=1.8cm] (c) at (-2.2,0.4) {client};
  \node[ellipse, draw, fill=lapp!50, minimum width=1.8cm] (s) at (2.2,0.4) {server};
  \node[draw, fill=lnet!60, rounded corners, align=center, font=\scriptsize\ttfamily] (lo) at (0,-0.6) {127.0.0.1\\localhost};
  \draw[{Stealth}-{Stealth}, thick, netblue] (c) to[bend right=25] (lo.west);
  \draw[{Stealth}-{Stealth}, thick, netblue] (lo.east) to[bend right=25] (s);
  \node (nic) at (5.4,0) {\icon[0.9cm]{switch}};
  \node[below=-1pt of nic, lbl] {verso la rete};
  \draw[black!50, line width=1.5pt] (h.east) -- (nic);
  \draw[netred, very thick] (3.95,0.25) -- (4.35,-0.25); \draw[netred, very thick] (3.95,-0.25) -- (4.35,0.25);
  \node[text=netred, lbl, anchor=north, text width=2cm] at (4.15,-0.35) {nulla raggiunge il cavo};
\end{tikzpicture}
\caption{Client e server sulla stessa macchina: è così che proverete i vostri programmi prima di avere un secondo host.}
\end{figure}

\subsection{Un nome non è un indirizzo}

Nessuno scrive \texttt{143.225.161.30}: si scrive \texttt{www.unina.it}, e i due designano lo stesso server Web.
\begin{itemize}
    \item Il \textbf{nome} è per le persone: leggibile, stabile, e sopravvive al cambio della macchina che c'è sotto.
    \item L'\textbf{indirizzo} è ciò su cui lavora il livello di rete: è ciò che serve all'instradamento per consegnare.
    \item Il \textbf{DNS} (\emph{Domain Name System}) traduce l'uno nell'altro. È un \textbf{database distribuito}
          (nessuna macchina lo contiene tutto) ed è a sua volta un'\textbf{applicazione} che gira sulla rete: lo
          apriremo tra poco.
\end{itemize}

% ══════════════════════════════════════════════════════════════
\section{I servizi di trasporto disponibili alle applicazioni}
% ══════════════════════════════════════════════════════════════

Il livello di trasporto, sotto l'applicazione, offre protocolli che realizzano su richiesta alcune
proprietà, dette \textbf{qualità del servizio} (QoS, \emph{Quality of Service}).

\begin{itemize}
    \item \textbf{Affidabilità} del trasferimento dati: i dati inviati da un estremo arrivano
          \textbf{corretti e completi} all'altro. Alcune applicazioni (file, posta, transazioni bancarie)
          non tollerano perdite; altre (audio e video) tollerano qualche perdita.
    \item \textbf{Throughput}: il ritmo, in bit/s, a cui il processo mittente può consegnare bit al
          processo ricevente. Alcune applicazioni multimediali hanno bisogno di un throughput minimo.
    \item \textbf{Temporizzazione} (\emph{timing}): ogni bit immesso nel socket arriva al socket del
          destinatario entro un certo tempo. Serve a telefonia, videoconferenza, giochi in tempo reale.
    \item \textbf{Sicurezza}: cifratura e decifratura dei messaggi, integrità, autenticazione.
\end{itemize}

Scegliere un protocollo di trasporto significa scegliere tra queste quattro cose. La domanda da farsi è:
\textbf{quale delle quattro serve davvero alla mia applicazione?}

\begin{table}[H]
\centering
\begin{tabular}{llll}
\toprule
\textbf{Applicazione} & \textbf{Perdite} & \textbf{Throughput} & \textbf{Sensibile al tempo} \\
\midrule
Trasferimento file & nessuna & elastico & no \\
Posta elettronica & nessuna & elastico & no \\
Documenti Web & nessuna & elastico, pochi kb/s & no \\
Telefonia su Internet & tollerate & da pochi kb/s a 1 Mb/s & sì, centinaia di ms \\
Videoconferenza & tollerate & da 10 kb/s a 5 Mb/s & sì, centinaia di ms \\
Audio/video registrato & tollerate & come sopra & sì, pochi secondi \\
Giochi interattivi & tollerate & da pochi kb/s a 10 kb/s & sì, centinaia di ms \\
\bottomrule
\end{tabular}
\caption{Che cosa chiedono al trasporto le applicazioni reali.}
\end{table}

\dfn{Throughput elastico}{
    Un'applicazione ha un throughput \textbf{elastico} se usa tutto il throughput che trova: di più è meglio, ma
    anche con meno funziona (più lentamente).
}

Internet offre alle applicazioni \textbf{due} protocolli di trasporto:

\begin{table}[H]
\centering
\begin{tabular}{p{2.3cm}p{6cm}p{6cm}}
\toprule
 & \textbf{TCP} (\emph{Transmission Control Protocol}) & \textbf{UDP} (\emph{User Datagram Protocol}) \\
\midrule
Connessione & \textbf{sì}: handshake tra i processi prima dei dati & \textbf{no}: si invia e basta \\
Affidabilità & trasferimento \textbf{affidabile} e ordinato & nessuna garanzia di consegna \\
Controlli & di flusso e di congestione & nessuno \\
Servizio & flusso di byte affidabile & datagram inaffidabile, ``minimale'' \\
Non offre & garanzie di tempo e di throughput minimo & garanzie di tempo e di throughput minimo \\
\bottomrule
\end{tabular}
\caption{I due servizi di trasporto di Internet.}
\end{table}

TCP rallenta il mittente quando la rete è congestionata: lo fa \textbf{per il bene della rete}, non della vostra
applicazione.

\info{UDP non è un TCP difettoso}{
    UDP è ciò che resta quando si \textbf{rifiuta di pagare} per le garanzie: niente handshake, niente conferme,
    niente rallentamenti imposti. Per un'applicazione che deve rispondere in fretta, o che gestisce da sé le
    perdite, è la scelta giusta.
}

\subsection{Che cosa Internet non garantisce}

Dei quattro servizi elencati all'inizio, i protocolli di trasporto di Internet ne forniscono \textbf{uno solo}:
l'affidabilità (con TCP).
\begin{itemize}
    \item \textbf{Nessuna garanzia di throughput né di tempo}: né TCP né UDP le offrono. Le applicazioni sensibili al
          tempo funzionano lo stesso perché sono \textbf{progettate per farcela}, non perché la rete le aiuti.
    \item \textbf{Nessuna sicurezza}: ciò che si spinge nel socket viaggia così come è stato scritto, password compresa.
\end{itemize}

\begin{figure}[H]
\centering
\begin{tikzpicture}[every node/.style={font=\footnotesize\sffamily}]
  \node[lay app, minimum width=3.2cm] (a) at (0,1.9) {Applicazione};
  \node[draw=netgreen!70!black, fill=netgreen!20, minimum width=3.2cm, minimum height=0.55cm] (tls) at (0,1.2) {TLS: cifra / decifra};
  \draw[very thick, netblue] (-1.9,0.8) -- (1.9,0.8) node[right, text=netblue] {socket};
  \node[lay tra, minimum width=3.2cm] (t) at (0,0.3) {TCP};
  \node[lay net, minimum width=3.2cm] at (0,-0.4) {IP \dots};
  \draw[decorate, decoration={brace, amplitude=5pt}] (-1.8,0.95) -- (-1.8,2.2) node[midway, left=6pt, lbl, text width=2.6cm, align=right] {livello applicativo};
  \node[lbl, text width=6cm, align=left, anchor=west] at (3.2,1.4) {l'applicazione cifra \textbf{prima} del socket e decifra \textbf{dopo}: TCP trasporta byte già illeggibili};
\end{tikzpicture}
\caption{TLS non è un terzo protocollo di trasporto: è un miglioramento di TCP, implementato nel livello applicativo.}
\end{figure}

% ══════════════════════════════════════════════════════════════
\section{I protocolli di livello applicativo}
% ══════════════════════════════════════════════════════════════

\dfn{Protocollo di livello applicativo}{
    Un \textbf{protocollo applicativo} definisce come i processi di un'applicazione, in esecuzione su end
    system diversi, si scambiano i messaggi. In particolare definisce:
    \begin{enumerate}
      \item i \textbf{tipi} di messaggi scambiati (per esempio richieste e risposte);
      \item la \textbf{sintassi} dei vari tipi di messaggio: quali campi ci sono e come sono delimitati;
      \item la \textbf{semantica} dei campi: il significato dell'informazione che contengono;
      \item le \textbf{regole} che stabiliscono quando e come un processo invia messaggi e risponde.
    \end{enumerate}
}

I protocolli possono essere \textbf{aperti} (le regole sono pubbliche, tipicamente in un RFC: HTTP,
SMTP, DNS) oppure \textbf{proprietari}, deliberatamente chiusi (come quello di Skype o di molte app di
messaggistica).

\subsection{Un protocollo non è l'applicazione}

Il Web è un'applicazione; HTTP ne è \textbf{un pezzo}.

\begin{figure}[H]
\centering
\begin{tikzpicture}[every node/.style={font=\footnotesize\sffamily}]
  \node[draw=primary, thick, rounded corners, minimum width=11.5cm, minimum height=2.5cm] (web) at (0,0) {};
  \node[anchor=north west, font=\small\bfseries\sffamily, text=primary] at (web.north west) {l'applicazione Web};
  \node[draw, fill=lapp!40, rounded corners, minimum width=2.2cm, minimum height=0.9cm] (b) at (-4.2,-0.2) {browser};
  \node[draw, fill=lapp!40, rounded corners, minimum width=2.2cm, minimum height=0.9cm] (s) at (4.2,-0.2) {server Web};
  \node[draw, fill=black!8, rounded corners, minimum width=2cm, minimum height=0.6cm] at (0,0.55) {formato HTML};
  \draw[{Stealth}-{Stealth}, very thick, netred] (b) -- node[below, text=netred, font=\footnotesize\bfseries\sffamily] {protocollo HTTP} (s);
  \node[lbl, text=netred] at (0,-0.95) {l'unica parte su cui i due capi devono mettersi d'accordo};
\end{tikzpicture}
\caption{Il Web ha bisogno di un formato per i documenti (HTML), di browser e di server: HTTP fissa soltanto formato e ordine dei messaggi tra browser e server.}
\end{figure}

Lo stesso vale per il video on demand: server, fatturazione, l'app sul telefono, e \textbf{DASH} per i messaggi.
Perché insistere sulla distinzione? Perché il protocollo è \textbf{l'unica parte su cui entrambi i capi devono
concordare}: tutto il resto ciascuna parte lo costruisce come vuole.

\subsection{Protocolli applicativi comuni}

Ecco alcuni protocolli applicativi di uso comune, con il trasporto che usano:

\begin{table}[H]
\centering
\begin{tabular}{lp{8.6cm}ll}
\toprule
\textbf{Protocollo} & \textbf{Descrizione} & \textbf{Porta} & \textbf{Trasporto} \\
\midrule
DHCP & \emph{Dynamic Host Configuration Protocol}: assegna gli indirizzi IP & 67, 68 & UDP \\
DNS & \emph{Domain Name System}: traduce i nomi in indirizzi IP & 53 & UDP, TCP \\
HTTP / HTTPS & \emph{HyperText Transfer Protocol (Secure)}: trasferisce pagine Web & 80 / 443 & TCP \\
SMTP / SMTPS & \emph{Simple Mail Transfer Protocol (Secure)}: invia la posta & 25 / 465 & TCP \\
SNMP & \emph{Simple Network Management Protocol}: gestisce i dispositivi di rete & 161 & UDP \\
Telnet / SSH & \emph{Teletype Network / Secure SHell}: riga di comando su host remoti & 23 / 22 & TCP \\
FTP / FTPS & \emph{File Transfer Protocol (Secure)}: trasferisce file & 21 / 990 & TCP \\
\bottomrule
\end{tabular}
\caption{Protocolli applicativi comuni e relative porte ben note.}
\end{table}

\info{Le versioni ``Secure''}{
    Molte applicazioni scambiano informazioni sensibili (credenziali, dati personali). Sulle reti pubbliche
    i messaggi vanno protetti: da qui le varianti sicure (HTTPS, SMTPS, FTPS, SSH) di quasi ogni protocollo.
    Telnet, per esempio, trasmette tutto in chiaro, password comprese; SSH offre lo stesso servizio (una
    riga di comando remota) con tutto cifrato. Su una rete pubblica, usare la variante in chiaro è un errore.
}

% ══════════════════════════════════════════════════════════════
\section{FTP e FTPS}
% ══════════════════════════════════════════════════════════════

\dfn{FTP}{
    L'\textbf{FTP} (\emph{File Transfer Protocol}) è uno dei protocolli più antichi di Internet (prima
    versione nel 1971) e serve a \textbf{trasferire file} tra host attraverso la rete, secondo il modello
    client-server.
}

\begin{itemize}
    \item Nella versione base il trasferimento avviene \textbf{in chiaro}: nome utente, password e file
          viaggiano leggibili da chiunque intercetti il traffico. Va quindi usato in reti locali o private.
    \item La versione sicura \textbf{FTPS} (\emph{FTP Secure}) protegge nome utente e password e cifra i contenuti.
    \item Sia FTP sia FTPS hanno due componenti: il \textbf{protocollo}, che specifica i comandi (elencare,
          prendere, cancellare file, \dots), e un'\textbf{applicazione} che lo implementa, con un lato client e
          un lato server.
\end{itemize}

\subsection{Due connessioni: controllo e dati}

A differenza dei protocolli visti finora, FTP \textbf{non} trasporta comandi e dati sulla stessa connessione.
\begin{itemize}
    \item La \textbf{connessione di controllo}, sulla porta \textbf{21}, si apre per prima e resta aperta per tutta
          la sessione: solo comandi e risposte.
    \item Ogni trasferimento (un \texttt{get}, un \texttt{put}, persino un \texttt{ls}) apre una \textbf{seconda
          connessione} solo per i dati, che viene chiusa alla fine del trasferimento.
    \item Tenere il controllo separato dai dati si chiama segnalazione \textbf{fuori banda} (\emph{out-of-band}).
\end{itemize}

\begin{figure}[H]
\centering
\begin{tikzpicture}[every node/.style={font=\footnotesize\sffamily}]
  \begin{scope}
    \node[font=\small\bfseries\sffamily] at (3,1.5) {Modalità attiva};
    \node (c) at (0,0) {\icon[1cm]{laptop}}; \node[below=-2pt of c] {client};
    \node (s) at (6,0) {\icon[0.75cm]{server}}; \node[below=-2pt of s] {server};
    \draw[-{Stealth}, very thick, netblue] (0.7,0.45) -- node[above] {controllo: il client apre (porta 21)} (5.6,0.45);
    \draw[-{Stealth}, very thick, netgreen!60!black] (5.6,-0.25) -- node[below] {dati: il \textbf{server} apre (dalla porta 20)} (0.7,-0.25);
    \node[draw=netred, fill=netred!8, rounded corners, text width=5.2cm, align=center] at (3,-1.75) {un client dietro NAT o firewall \textbf{non} è raggiungibile dall'esterno: la connessione dati fallisce};
  \end{scope}
  \begin{scope}[xshift=8.2cm]
    \node[font=\small\bfseries\sffamily] at (3,1.5) {Modalità passiva};
    \node (c) at (0,0) {\icon[1cm]{laptop}}; \node[below=-2pt of c] {client};
    \node (s) at (6,0) {\icon[0.75cm]{server}}; \node[below=-2pt of s] {server};
    \draw[-{Stealth}, very thick, netblue] (0.7,0.45) -- node[above] {controllo: il client apre (porta 21)} (5.6,0.45);
    \draw[-{Stealth}, very thick, netgreen!60!black] (0.7,-0.25) -- node[below] {dati: apre \textbf{ancora il client}} (5.6,-0.25);
    \node[draw=netgreen!70!black, fill=netgreen!8, rounded corners, text width=5.2cm, align=center] at (3,-1.75) {esiste proprio per questo: il client apre entrambe le connessioni};
  \end{scope}
\end{tikzpicture}
\caption{Chi apre la connessione dati? In modalità attiva il server si ricollega al client; in modalità passiva è il client ad aprire entrambe.}
\end{figure}

\subsection{Una sessione FTP}

Sul server un \textbf{demone} attende in background le connessioni; sul client si usa il programma \texttt{ftp}.

\begin{lstlisting}[language=bash]
$ sudo apt install vsftpd          # lato server, una volta sola
$ systemctl status vsftpd          # il demone e' in esecuzione?

$ ftp 143.225.161.30               # lato client
Name: student
Password:
230 Login successful.
ftp> ls                            # da qui in poi, il prompt ftp>
ftp> get report.pdf                # dal server al client
226 Transfer complete.
ftp> bye                           # bye o quit, non exit
\end{lstlisting}

I codici numerici (230, 226) sono le \textbf{risposte} del server sulla connessione di controllo: 230 conferma
l'accesso, 226 la fine del trasferimento.

\info{Il demone}{
    Il lato server è implementato come \textbf{demone} (\emph{daemon}): un programma che gira in background
    e resta in attesa che i client si connettano. Quasi ogni server che scriverete in questo corso ha questa
    forma: un demone, in attesa.
}

\begin{table}[H]
\centering
\begin{tabular}{ll}
\toprule
\textbf{Comando} & \textbf{Effetto} \\
\midrule
\texttt{help} & elenca i comandi FTP disponibili \\
\texttt{cd} / \texttt{lcd} & cambia directory sulla macchina remota / locale \\
\texttt{ls} & mostra file e directory nella directory remota corrente \\
\texttt{pwd} & stampa la directory di lavoro remota \\
\texttt{mkdir} / \texttt{rmdir} & crea / rimuove una directory remota \\
\texttt{delete} & cancella un file nella directory remota \\
\texttt{get} & copia un file dal server alla macchina locale (\emph{download}) \\
\texttt{put} & copia un file dalla macchina locale al server (\emph{upload}) \\
\bottomrule
\end{tabular}
\caption{I comandi FTP principali.}
\end{table}

Due comandi, \texttt{get} e \texttt{put}, sono la ragione per cui il protocollo esiste; il resto è navigazione.
Attenzione a \texttt{lcd}: è l'unico che agisce sulla \textbf{vostra} macchina.

% ══════════════════════════════════════════════════════════════
\section{Riepilogo}
% ══════════════════════════════════════════════════════════════

\riepilogo{
\begin{itemize}
    \item Un'applicazione di rete è fatta di programmi su \textbf{end system diversi}; il nucleo della rete
          non esegue applicazioni.
    \item Architetture: \textbf{client-server} (server sempre attivo, indirizzo noto) e \textbf{P2P} (peer
          intermittenti, scalabile); nella pratica spesso \textbf{ibride}.
    \item Comunicano \textbf{processi}, attraverso i \textbf{socket}: l'interfaccia tra applicazione e trasporto.
    \item Un processo si raggiunge con \textbf{indirizzo IP} (host, livello di rete) + \textbf{numero di porta}
          (processo, livello di trasporto).
    \item Un indirizzo IP appartiene a un'\textbf{interfaccia}; \texttt{127.0.0.1} indica l'host stesso; il DNS
          traduce i nomi (per le persone) in indirizzi (per la rete).
    \item Le applicazioni chiedono al trasporto affidabilità, throughput, tempi, sicurezza; Internet offre
          \textbf{TCP} (connessione, affidabile) e \textbf{UDP} (senza connessione, minimale), e garantisce solo
          l'affidabilità: niente throughput, tempi o sicurezza (TLS si aggiunge sopra TCP, a livello applicativo).
    \item Un protocollo applicativo definisce tipi, sintassi, semantica dei messaggi e regole di scambio, ed è
          solo \textbf{un pezzo} dell'applicazione.
    \item \textbf{FTP} usa due connessioni: controllo (porta 21, per tutta la sessione) e dati (una per
          trasferimento); modalità attiva e passiva.
\end{itemize}
}
